\section{Background}
\label{sec:background}

\subsection{Preliminaries}
\label{sec:preliminaries}

Conjunctive queries represent information needs as relations that must
hold jointly. A query over binary relations has the form
\begin{equation}
 Q(x)=\exists\bar y\;\bigwedge_{j=1}^{m}r_j(t_j,t'_j),
 \label{eq:cq}
\end{equation}
where each term is a constant or variable, $x$ is the answer variable,
and $\bar y$ contains the remaining
variables~\cite{chandra1977conjunctive}. An answer is obtained by assigning
values to the variables so that every atom matches a stored fact.
Occurrences of the same variable must receive the same value.

A \emph{witness} is a set of facts supporting such an assignment: the
query yields the answer using those facts
alone~\cite{buneman2001why}. A witness contains at most $m$ distinct facts, and different witnesses can support the same answer.
Provenance records these alternative
derivations~\cite{green2007provenance}. This perspective gives retrieval
a structural target: finding facts that jointly satisfy a query, rather
than selecting them solely by their individual relevance.

\subsection{Related Work}
\label{sec:related}

\paragraph{Memory architectures for language-model agents.}
Existing systems address long interaction histories through virtual
context management~\cite{packer2023memgpt,lee2024readagent}, memory streams
with reflection or forgetting~\cite{park2023generative,zhong2024memorybank},
and feedback accumulated across
episodes~\cite{shinn2023reflexion,flexa2026bridging}. Other architectures
maintain extracted facts~\cite{chhikara2025mem0}, linked
notes~\cite{xu2025amem}, temporal
graphs~\cite{rasmussen2025zep,anokhin2025arigraph}, hierarchical
stores~\cite{kang2025memoryos,zhang2025gmemory}, or compressed and enriched
records~\cite{fang2026lightmem,liu2026simplemem,salama2025meminsight}.
These designs illustrate the range of representations available to a
memory system. \sys focuses on how a question's relational requirements
can guide retrieval from such a representation. Its stored propositions
serve both as readable memory records and as graph facts with source
provenance.

\paragraph{Graph-based and iterative retrieval.}
Graph and hierarchical retrieval use structure to recover information
beyond direct similarity to the question. Examples include community
summaries~\cite{edge2024graphrag}, recursive summary
trees~\cite{sarthi2024raptor}, and relevance propagation over open
knowledge graphs~\cite{gutierrez2024hipporag,gutierrez2025hipporag2}.
Iterative approaches instead interleave retrieval with
reasoning~\cite{trivedi2023ircot}, guide graph exploration with a language
model~\cite{sun2024tog}, or adapt retrieval to the
input~\cite{asai2024selfrag,jeong2024adaptiverag}. In agent memory,
MemAgent learns memory behavior through reinforcement
learning~\cite{yu2026memagent}, while GAM uses a researcher that searches
stored pages, integrates evidence, and reflects on whether further
search is needed~\cite{yan2025gam}. \sys also revises retrieval in
response to evidence. Its distinguishing design choice is to expose a
conjunctive plan whose bindings and supporting facts are computed by
an executor, with failed execution conditions returned to the planner.

\paragraph{Temporal memory.}
Time can be represented as a retrieval preference or as an attribute of
stored information. Generative Agents use access-dependent
recency~\cite{park2023generative}; MemoryBank incorporates forgetting over
elapsed time~\cite{zhong2024memorybank}; Zep represents temporal
information in graph relations~\cite{rasmussen2025zep}; and Zero-Mem
organizes evidence using temporal structure~\cite{xiao2026zeromem}.
Recency alone does not express every temporal information need: a
question may refer to an early event or a named historical period.
\sys lets the planner choose the reference against which temporal
proximity is measured. This reference changes ranking rather than
imposing a hard temporal constraint on graph execution.

\paragraph{Provenance, attribution, and verification.}
Database provenance characterizes the facts contributing to a query
answer~\cite{buneman2001why,green2007provenance}. In language models,
attribution associates generated claims with source
passages~\cite{gao2023alce}, while verification can use additional
questions to identify and revise unsupported
claims~\cite{dhuliawala2024cove}. Source-linked retrieval is also relevant
to agent-memory designs such as Zero-Mem~\cite{xiao2026zeromem}.
\sys uses provenance within retrieval: a candidate witness identifies
the source utterances supplied to verification and influences context
selection. This makes the supporting derivation inspectable, while
leaving the verifier and final reader subject to model error.
